\documentclass[a4paper,11pt]{article}

% --- Pacchetti Configurazione ---
\usepackage[utf8]{inputenc}
\usepackage[italian]{babel}
\usepackage[T1]{fontenc}
\usepackage{geometry}
\geometry{margin=2.5cm}
\usepackage{booktabs}
\usepackage{xcolor}
\usepackage{titlesec}
\usepackage{enumitem}
\usepackage{amsmath}
\usepackage{amssymb}
\usepackage{hyperref}
\usepackage{array}

% --- Definizione Colori e Stile ---
\definecolor{primaryblue}{rgb}{0.1, 0.3, 0.6}
\definecolor{grayline}{gray}{0.8}

\titleformat{\section}{\large\bfseries\color{primaryblue}}{\thesection}{1em}{}[\titlerule]
\titleformat{\subsection}{\normalsize\bfseries\color{primaryblue}}{\thesubsection}{1em}{}

\hypersetup{
	colorlinks=true,
	linkcolor=primaryblue,
	urlcolor=cyan,
	pdftitle={Report LabOS 2.0 - BicoccaLab v7},
}

% --- Metadati ---
\title{\textbf{Report di Stato Avanzamento: LabOS 2.0}\\ \large (BicoccaLab v7)}
\author{Dipartimento di Automazione e Robotica}
\date{20 aprile 2026}

\begin{document}
	
	\maketitle
	
	\section{Introduzione e Visione Strategica}
	Il progetto \textbf{LabOS 2.0} consolida una piattaforma modulare progettata per eliminare il \textit{"Time-to-Experiment bottleneck"}. Con un \textbf{Technology Readiness Level (TRL) 6/7}, il sistema transita verso un modello di \textbf{Self-Driving Lab (SDL)}.
	
	I quattro pilastri strategici identificati sono:
	\begin{itemize}[label=$\checkmark$]
		\item \textbf{Coordinamento eterogeneo:} Integrazione multi-vendor unificata.
		\item \textbf{Adattabilità:} Riconfigurazione rapida dei protocolli (\textit{No-Code}).
		\item \textbf{Human-in-the-loop:} Sinergia uomo-macchina per task complessi.
		\item \textbf{Ottimizzazione Closed-loop:} Generazione dati standardizzati per AI.
	\end{itemize}
	
	\section{Infrastruttura Operativa e Workstation}
	L'architettura garantisce \textbf{scalabilità orizzontale}, permettendo l'aggiunta di moduli senza alterare la logica centrale.
	
	
	
	\begin{table}[h!]
		\centering
		\small
		\caption{Componenti Standard e Impatto Strategico}
		\vspace{2mm}
		\begin{tabular}{@{}p{3cm}p{4cm}p{7cm}@{}}
			\toprule
			\textbf{Componente} & \textbf{Ruolo Operativo} & \textbf{Impatto Strategico (\textit{So What?})} \\ \midrule
			\textbf{Opentrons Flex} & Pipettaggio e preparazione. & HAL e API Python riducono i tempi di ri-validazione del 40\%. \\ \addlinespace
			\textbf{Tecan M200 Pro} & Analisi spettroscopica. & Integrazione SiLA2 elimina i silos informativi (XML $\to$ Strutturato). \\ \addlinespace
			\textbf{GoFaGo} & Logistica autonoma. & Cinematica omnidirezionale ottimizza il \textit{footprint} di laboratorio. \\ \addlinespace
			\textbf{Manual Station} & Intervento assistito. & Garantisce resilienza in caso di variabilità estrema dei materiali. \\ \bottomrule
		\end{tabular}
	\end{table}
	
	\section{Standardizzazione: Il Motore SiLA2 e HAL}
	L'adozione dello standard \textbf{SiLA2} (basato su gRPC) garantisce comunicazioni binarie efficienti. L'\textbf{Hardware Abstraction Layer (HAL)} risolve a \textit{runtime} le \textit{logical labware} in \textit{physical slots}, prevenendo errori di \textit{hard-coding}.
	
	\subsection{Flusso di Esecuzione del Workflow}
	\begin{enumerate}
		\item \textbf{Definizione JSON:} Formato dichiarativo leggibile dall'uomo e dalla macchina.
		\item \textbf{Risoluzione HAL:} Mapping della ricetta sulla configurazione fisica attuale.
		\item \textbf{Esecuzione:} Generazione dinamica di protocolli Python e comandi SiLA2.
		\item \textbf{Standardizzazione:} Consolidamento dei risultati in formato \textbf{AnIML}.
	\end{enumerate}
	
	\section{Validazione Sperimentale: Sintesi Idrogeli}
	I test di stress hanno definito l'inviluppo di tolleranza per il docking autonomo:
	\begin{itemize}
		\item \textbf{Task Tolerance Laterale (X):} $\pm 4$ cm (limite cinematico del braccio).
		\item \textbf{Task Tolerance Longitudinale (Y):} da $-5$ a $+20$ cm.
	\end{itemize}
	
	\textbf{Criticità identificate:} Il \textbf{Tecan Carrier} presenta il tasso di successo più basso (80\%) a causa della struttura a sbalzo (\textit{cantilevered tray}). È stato rilevato un \textit{drift} di calibrazione della camera di $1.8$ mm e un effetto \textit{"forward-push"} indotto dalla pinza RG6.
	
	\section{Analisi del Rischio e Azioni Correttive}
	\begin{table}[h!]
		\centering
		\begin{tabular}{@{}lll@{}}
			\toprule
			\textbf{Rischio} & \textbf{Livello} & \textbf{Strategia di Mitigazione} \\ \midrule
			Tecnologico & Medio-Basso & Manutenzione preventiva e monitoraggio KPI. \\
			Adozione & Medio & Formazione tramite interfaccia \textit{No-Code}. \\
			Scalabilità & Basso & Architettura SiLA2 nativamente estensibile. \\ \bottomrule
		\end{tabular}
	\end{table}
	
	\textbf{Azioni Correttive Prioritarie:}
	\begin{itemize}
		\item Sostituzione pinza RG6 con modello a \textbf{cinematica parallela}.
		\item Passaggio a \textbf{marker ArUco in alluminio} per prevenire deformazioni.
		\item Stress-test 24/7 per monitoraggio derive termiche.
	\end{itemize}
	
	\section{Roadmap Operativa (90 Giorni)}
	\begin{description}
		\item[Mese 1: Hardening] Rilascio pacchetto di qualifica e criteri di accettazione.
		\item[Mese 2: Monitoring] Report dettagliato su Throughput e \textit{Intervention Rate}.
		\item[Mese 3: Readiness] Dimostrazione \textit{end-to-end} e transizione operativa finale.
	\end{description}
	
\end{document}